% Retrospectiva do desenvolvimento; não representa uma apresentação realizada no passado.
% Este PDF não inclui resultados do teste reservado.
\documentclass[aspectratio=169,11pt]{beamer}
\usepackage[T1]{fontenc}
\usepackage[utf8]{inputenc}
\usepackage[brazil]{babel}
\usepackage{lmodern,booktabs,amsmath,graphicx}
\definecolor{azul}{HTML}{172D3C}
\definecolor{verde}{HTML}{087F8C}
\setbeamercolor{normal text}{fg=azul,bg=white}
\setbeamercolor{frametitle}{fg=azul}
\setbeamercolor{structure}{fg=verde}
\setbeamercolor{block title}{fg=white,bg=verde}
\setbeamercolor{block body}{bg=verde!6}
\setbeamerfont{frametitle}{size=\Large,series=\bfseries}
\setbeamertemplate{navigation symbols}{}
\setbeamertemplate{itemize items}[circle]
\setbeamertemplate{footline}{\hspace{5mm}\scriptsize ENEM 2023 | Unifor\hfill\insertframenumber/\inserttotalframenumber\hspace{5mm}\vspace{3mm}}
\setbeamersize{text margin left=8mm,text margin right=8mm}
\graphicspath{{figuras/}}
\newcommand{\figura}[1]{\begin{center}\includegraphics[width=.96\linewidth,height=5.35cm,keepaspectratio]{#1}\end{center}}
\newcommand{\fonte}[1]{\par\vspace{1mm}{\tiny #1\par}}
\newcommand{\duas}[2]{\begin{columns}[T,onlytextwidth]\begin{column}{.48\textwidth}#1\end{column}\begin{column}{.48\textwidth}#2\end{column}\end{columns}}
\title{Comparação nacional de modelos\\ENEM 2023}
\author{Joey Alan de Freitas Solis \and Paulo Henrico}
\institute{Universidade de Fortaleza (Unifor)\\Professor Tulio}
\date{}

\begin{document}
\begin{frame}\titlepage\centering\textcolor{verde}{Acompanhamento: retrospectiva dos experimentos no desenvolvimento}\end{frame}
\begin{frame}{Base, alvo e participantes}
\duas{\begin{tabular}{lr}\toprule Base nacional & 3.933.955\\Colunas originais & 76\\Sem alvo observado & 1.241.528\\Elegíveis & 2.692.427\\Desenvolvimento & 2.153.941\\Teste reservado & 538.486\\\bottomrule\end{tabular}}{\begin{itemize}\item Alvo: NU\_NOTA\_CN.
\item 16.547 zeros mantidos; ausência não é zero.
\item Fonte real: microdados do Inep, ENEM 2023.
\item Este acompanhamento utiliza apenas resultados do desenvolvimento.\end{itemize}}
\end{frame}
\begin{frame}{Features, limpeza e representação}
\begin{itemize}\item 40 entradas: moradores, perfil demográfico/escolar e questionário.
\item ID somente para controle; outras notas, presença e respostas fora de X.
\item Ingestão valida tipos, ausências, ano, IDs e contagens, preservando originais.
\item Q005: mediana, indicador quando aplicável e padronização.
\item Categorias: ausência explícita + one-hot; 235 dimensões originais.
\item Estatísticas de preparação ajustadas no treino de cada dobra.\end{itemize}
\fonte{Fontes de implementação: configs/ciencias\_natureza.toml; cleaning.py; preprocessing.py.}
\end{frame}
\begin{frame}{Distribuição do alvo no desenvolvimento}
\figura{01_distribuicao.png}
\fonte{Inep, ENEM 2023. Notas extremas são menos frequentes; zero é resultado observado.}
\end{frame}
\begin{frame}{Protocolo e validação cruzada}
\begin{block}{Partições comuns}80\% desenvolvimento e 20\% teste, random\_state=42; três dobras no desenvolvimento, persistidas em dados/splits/.\end{block}
\centering\begin{tabular}{lcc}\toprule Rodada & Treino & Validação\\\midrule 1 & B+C & A\\2 & A+C & B\\3 & A+B & C\\\bottomrule\end{tabular}
\vspace{3mm}\begin{itemize}\item RMSE principal; MAE e R\textsuperscript{2} complementares.
\item Comparar média e DP das mesmas dobras.
\item Não escolher solução com resultados do teste.\end{itemize}
\end{frame}
\begin{frame}{Resultados da comparação inicial}
\figura{18_comparacao_completa.png}
\fonte{Baseline: RMSE 87,921. Stacking: 76,175. Todos possuem três dobras completas; nenhuma métrica de teste nesta figura.}
\end{frame}
\begin{frame}{Investigações seguintes e questões de pesquisa}
\begin{itemize}\item Criar três atributos: recursos por morador, acesso digital e escolaridade.
\item Seleção, PCA e ponderação das faixas raras melhoram a previsão?
\item Quanto muda o erro ao retirar questionário ou UF?
\item Comparar parâmetros dos melhores regressores e curvas progressivas.
\item Congelar a solução e somente depois avaliar o teste final.
\end{itemize}
\fonte{Retrospectiva: essas investigações foram posteriormente executadas e estão documentadas no artigo e no pitch. Este PDF não comprova apresentação ou submissão anterior.}
\end{frame}
\end{document}
